\documentclass[a4paper,fontset=none]{ctexart}

\usepackage{indentfirst}
\setlength{\parindent}{0em}
\usepackage{autobreak}
\allowdisplaybreaks
\usepackage{parskip}
\usepackage{geometry}
\geometry{top=3cm,bottom=3cm,left=2cm,right=2cm}
\usepackage{anyfontsize}

\usepackage{fontspec}
\setmainfont{STIX Two Text}
\usepackage{unicode-math}
\unimathsetup{mathrm=sym,mathbf=sym,mathsf=sym,bold-style=ISO}
\setmathfont{STIX Two Math}
\usepackage{physics2}
\usephysicsmodule{ab,op.legacy}
\usepackage{fixdif,derivative}
\newcommand{\um}{\upmu\mathrm{m}}
\newcommand{\ang}{\text{\r{A}}}
\AtBeginDocument{
    \let\uglyepsilon\epsilon
    \renewcommand\epsilon{\varepsilon}
}

\setCJKmainfont{Source Han Serif CN}[BoldFont=Source Han Sans CN Medium,ItalicFont=FZKai-Z03]

\title{\textbf{星际介质天文学作业}}
\author{张植竣}
\date{2024年3月22日}

\begin{document}

\maketitle

\begin{enumerate}
    \item[22.1 (a)]
    因为$\lambda = 10\,\um \ll a = 0.1\,\um$，可以认为这里的尘埃可以用电偶极近似，由(22.19)式得：
    \[
        C_{\mathrm{abs}}
        = 18\pi \frac{\epsilon_2}{(\epsilon_1 + 2)^2 + \epsilon_2^2} \cdot \frac{V}{\lambda}
        = 18\pi \times \frac{2.575}{(0.996 + 2)^2 + 2.575^2} \times \frac{(4\pi/3)a^3}{\lambda} = 3.908 \times 10^{-3}\,\um^2
    \]
    则有：
    \[
        Q_{\mathrm{abs}} = \frac{C_{\mathrm{abs}}}{\pi a^2} = 0.124
    \]
    且：
    \[
        \frac{C_{\mathrm{abs}}}{V} = \frac{C_{\mathrm{abs}}}{(4\pi/3)a^3} = 0.933\,\um^{-1}
    \]

    \medskip
    \item[22.1 (b)]
    首先使用(22.20)式计算$C_{\mathrm{sca}}$，注意这里的$\epsilon$是复数：
    \[
        C_{\mathrm{sca}}
        = 24\pi^3 \ab|\frac{\epsilon - 1}{\epsilon + 2}|^2 \frac{V^2}{\lambda^4}
        = 24\pi^3 \times \ab|\frac{-0.004 + 2.575i}{2.996 + 2.575i}|^2 \times \frac{\ab[(4\pi/3)a^3]^2}{\lambda^4} = 5.547 \times 10^{-7}\,\um^2
    \]
    则有：
    \[
        Q_{\mathrm{sca}} = \frac{C_{\mathrm{sca}}}{\pi a^2} = 1.766 \times 10^{-5}
    \]

    \bigskip
    \item[22.2 (a)]
    由(22.12)和(22.14)式，当尘埃的尺度远小于电磁波的波长时，吸收截面为：
    \[
        C_{\mathrm{abs},j} = \frac{4\pi\omega}{c}\Im(\alpha_{jj}) = \frac{4\pi\omega}{c}\Im\ab[\frac{V}{4\pi} \cdot \frac{\epsilon - 1}{(\epsilon - 1)L_j + 1}]
    \]
    首先计算$\mathrm{Im}(\alpha_{jj})$，需要将$\alpha_{jj}$的实部与虚部分离：
    \begin{align*}
        \alpha_{jj}
        &= \frac{V}{4\pi} \cdot \frac{\epsilon - 1}{(\epsilon - 1)L_j + 1} \\
        &= \frac{V}{4\pi} \cdot \frac{(\epsilon_1 - 1) + i\epsilon_2}{1 + (\epsilon_1 - 1)L_j + i\epsilon_2 L_j} \\
        &= \frac{V}{4\pi} \cdot \frac{[(\epsilon_1 - 1) + i\epsilon_2] \cdot [1 + (\epsilon_1 - 1)L_j - i\epsilon_2 L_j]}{[1 + (\epsilon_1 - 1)L_j]^2 + (\epsilon_2 L_j)^2} \\
        &= \frac{V}{4\pi} \cdot \frac{(\epsilon_1 - 1) \cdot [1 + (\epsilon_1 - 1)L_j] + \epsilon_2^2 L_j + i\epsilon_2}{[1 + (\epsilon_1 - 1)L_j]^2 + (\epsilon_2 L_j)^2}
    \end{align*}
    则其虚部为：
    \[
        \mathrm{Im}(\alpha_{jj}) = \frac{V}{4\pi} \cdot \frac{\epsilon_2}{[1 + (\epsilon_1 - 1)L_j]^2 + (\epsilon_2 L_j)^2}
    \]
    于是有：
    \[
        C_{\mathrm{abs},j} = \frac{4\pi\omega}{c}\mathrm{Im}(\alpha_{jj}) = \frac{\omega V}{c} \cdot \frac{\epsilon_2}{[1 + (\epsilon_1 - 1)L_j]^2 + (\epsilon_2 L_j)^2}
    \]
    又因为$\omega = 2\pi\nu$、$\lambda\nu = c$，则$\dfrac{\omega}{c} = \dfrac{2\pi}{\lambda}$，故：
    \[
        C_{\mathrm{abs},j} = \frac{2\pi V}{\lambda} \cdot \frac{\epsilon_2}{[1 + (\epsilon_1 - 1)L_j]^2 + (\epsilon_2 L_j)^2}
    \]

    \medskip
    \item[22.2 (b)]
    由题意：
    \begin{align*}
        C_{\mathrm{abs},a} - C_{\mathrm{abs},b}
        &= \frac{2\pi V}{\lambda} \times \ab\{\frac{\epsilon_2}{[1 + (\epsilon_1 - 1)L_a]^2 + (\epsilon_2 L_a)^2} - \frac{\epsilon_2}{[1 + (\epsilon_1 - 1)L_b]^2 + (\epsilon_2 L_b)^2}\} \\
        &= \frac{2\pi V\epsilon_2}{\lambda} \times \frac{\ab\{[1 + (\epsilon_1 - 1)L_b]^2 + (\epsilon_2 L_b)^2\} - \ab\{[1 + (\epsilon_1 - 1)L_a]^2 + (\epsilon_2 L_a)^2\}}{\ab\{[1 + (\epsilon_1 - 1)L_a]^2 + (\epsilon_2 L_a)^2\} \cdot \ab\{[1 + (\epsilon_1 - 1)L_b]^2 + (\epsilon_2 L_b)^2\}} \\
        &= \frac{2\pi V\epsilon_2}{\lambda} \times \frac{\ab\{[1 + (\epsilon_1 - 1)L_b]^2 - [1 + (\epsilon_1 - 1)L_a]^2\} + \ab\{(\epsilon_2 L_b)^2 - (\epsilon_2 L_a)^2\}}{\ab\{[1 + (\epsilon_1 - 1)L_a]^2 + (\epsilon_2 L_a)^2\} \cdot \ab\{[1 + (\epsilon_1 - 1)L_b]^2 + (\epsilon_2 L_b)^2\}} \\
        &= \frac{2\pi V\epsilon_2}{\lambda} \times \frac{[2 + (\epsilon_1 - 1)L_b + (\epsilon_1 - 1)L_a] \cdot [(\epsilon_1 - 1)(L_b - L_a)] + \epsilon_2^2 (L_b + L_a) \cdot (L_b - L_a)}{\ab\{[1 + (\epsilon_1 - 1)L_a]^2 + (\epsilon_2 L_a)^2\} \cdot \ab\{[1 + (\epsilon_1 - 1)L_b]^2 + (\epsilon_2 L_b)^2\}} \\
        &= \frac{2\pi V\epsilon_2}{\lambda} (L_b - L_a) \times \frac{[2 + (\epsilon_1 - 1) (L_b + L_a)] \cdot (\epsilon_1 - 1) + \epsilon_2^2 (L_b + L_a)}{\ab\{[1 + (\epsilon_1 - 1)L_a]^2 + (\epsilon_2 L_a)^2\} \cdot \ab\{[1 + (\epsilon_1 - 1)L_b]^2 + (\epsilon_2 L_b)^2\}} \\
        &= \frac{2\pi V\epsilon_2}{\lambda} (L_b - L_a) \times \frac{(\epsilon_1 - 1)^2 (L_b + L_a) + 2(\epsilon_1 - 1) + \epsilon_2^2 (L_b + L_a)}{\ab\{[1 + (\epsilon_1 - 1)L_a]^2 + (\epsilon_2 L_a)^2\} \cdot \ab\{[1 + (\epsilon_1 - 1)L_b]^2 + (\epsilon_2 L_b)^2\}}
    \end{align*}

    \bigskip
    \item[22.3 (a)]
    首先计算$e$：
    \[
        e = \sqrt{1 - \ab(\frac{b}{a})^2} = \frac{\sqrt{5}}{2}
    \]
    因为$a < b$，由(22.16)式得：
    \[
        L_a = \frac{1+e^2}{e^2} \ab(1 - \frac{1}{e}\arctan{e}) = 0.446
    \]
    则：
    \[
        L_b = \frac{1}{2}(1 - L_a) = 0.277
    \]

    \medskip
    \item[22.3 (b)]
    由(22.2)题得：
    \[
        C_{\mathrm{abs},j} = \frac{2\pi V}{\lambda} \cdot \frac{\epsilon_2}{[1 + (\epsilon_1 - 1)L_j]^2 + (\epsilon_2 L_j)^2}
    \]
    分别代入$L_a$、$L_b$得：
    \[
        \frac{C_{\mathrm{abs},a}}{V} = \frac{2\pi}{\lambda} \cdot \frac{\epsilon_2}{[1 + (\epsilon_1 - 1)L_a]^2 + (\epsilon_2 L_a)^2} = 0.699\,\um^{-1}
    \]
    \[
        \frac{C_{\mathrm{abs},b}}{V} = \frac{2\pi}{\lambda} \cdot \frac{\epsilon_2}{[1 + (\epsilon_1 - 1)L_b]^2 + (\epsilon_2 L_b)^2} = 1.074\,\um^{-1}
    \]
    则：
    \[
        \frac{C_{\mathrm{abs},b} - C_{\mathrm{abs},a}}{V} = 1.074\,\um^{-1} - 0.699\,\um^{-1} = 0.375\,\um^{-1}
    \]

    \medskip
    \item[22.3 (c)]
    如果这些椭球形尘埃是随机分布的，平均的吸收截面应当是按比例分配：
    \[
        \frac{\langle C_{\mathrm{abs}} \rangle}{V} = \frac{1}{3} \times \frac{C_{\mathrm{abs},a}}{V} + \frac{2}{3} \times\frac{C_{\mathrm{abs},b}}{V} = 0.949\,\um^{-1}
    \]
    由(22.1)题得同类型的$a = 0.1\,\um$球形尘埃在$\lambda = 10\,\um$的吸收截面为$\dfrac{C_{\mathrm{abs}}}{V} = 0.933\,\um^{-1}$，则$a = 0.1\,\um$、$b = 0.15\,\um$的椭球形尘埃的吸收截面比同类型的$a = 0.1\,\um$球形尘埃要偏大约$1.7\%$。

    \bigskip
    \item[22.4 ]
    由题意，波长为$\lambda = 450\,\um$时，介电常数为：
    \[
        \epsilon \approx 11.6 + 3.0 \times \ab(\frac{100\,\um}{450\,\um})i = 11.6 + 0.667i
    \]
    即：
    \[
        \epsilon_1 = 11.6,\quad \epsilon_2 = 0.667
    \]
    与(22.3)题相同，对于这种椭球形的尘埃粒子有：
    \[
        L_a = 0.446,\quad L_b = 0.277
    \]
    由(22.2)题得：
    \[
        C_{\mathrm{abs},j} = \frac{2\pi V}{\lambda} \cdot \frac{\epsilon_2}{[1 + (\epsilon_1 - 1)L_j]^2 + (\epsilon_2 L_j)^2}
    \]
    分别代入$L_a$、$L_b$得：
    \[
        \frac{C_{\mathrm{abs},a}}{V} = \frac{2\pi}{\lambda} \cdot \frac{\epsilon_2}{[1 + (\epsilon_1 - 1)L_a]^2 + (\epsilon_2 L_a)^2} = 2.831\,\um^{-1}
    \]
    \[
        \frac{C_{\mathrm{abs},b}}{V} = \frac{2\pi}{\lambda} \cdot \frac{\epsilon_2}{[1 + (\epsilon_1 - 1)L_b]^2 + (\epsilon_2 L_b)^2} = 5.993\,\um^{-1}
    \]
    则：
    \[
        \frac{C_{\mathrm{abs},b}}{C_{\mathrm{abs},a}} = \frac{C_{\mathrm{abs},b}}{V} \div \frac{C_{\mathrm{abs},a}}{V} = 2.117
    \]
    因为$\hat{a} \perp \mathbf{B}_0$，则垂直于$\mathbf{B}_0$的方向，即垂直于$\hat{a}$的方向，热辐射的偏振度为：
    \begin{align*}
        p
        &= \frac{I_{\mathrm{abs},b} - I_{\mathrm{abs},a}}{I_{\mathrm{abs},b} + I_{\mathrm{abs},a}} \\
        &= \frac{C_{\mathrm{abs},b} - C_{\mathrm{abs},a}}{C_{\mathrm{abs},b} + C_{\mathrm{abs},a}} \\
        &= \frac{(C_{\mathrm{abs},b} \div C_{\mathrm{abs},a}) - 1}{(C_{\mathrm{abs},b} \div C_{\mathrm{abs},a}) + 1} \\
        &= \frac{2.117 - 1}{2.117 + 1} \\
        &= 0.358
    \end{align*}

\end{enumerate}

\end{document}
